\chapter{Nguồn dữ liệu và chi tiết tái lập}
Các đường dẫn tính từ thư mục gốc \texttt{Thesis\_26}.
Tập bằng chứng ngày 12/11/2026 lưu tại \path{docs/thesis/evidence/snapshot.json}.
Tệp \path{input_hashes.json} cùng thư mục ghi mã SHA-256 của các nguồn đã đối chiếu.

\section{Cấu hình và bản ghi thực nghiệm}
Bản ghi đánh giá vòng kín được lưu tại
\begin{quote}\small
\path{experiments/evaluation/results_mem_2/step5_test_20261109-093744}
\end{quote}
Mỗi điều kiện có bản kê \path{run_manifest.json}, phân hoạch đã khóa
\path{split.lock.json} và tóm tắt từng lượt trong \path{episodes/*/summary.json}.
Phép tổng hợp đối chiếu trạng thái hoàn tất và SHA-256 với \path{status.json}.
Bản kê lưu thiết lập điều khiển, danh sách trạng thái khởi tạo, phiên bản mã và mô hình.

Bản ghi huấn luyện được lưu tại
\begin{quote}\small
\path{experiments/training/results_mem_2/mem-adapter-20261105-083543}
\end{quote}
Nguồn dùng để đối chiếu gồm \path{run_config.json}, \path{metrics.jsonl},
\path{checkpoint-003000/progress.json} và bản xuất trong \path{pretrained_model}.
Danh mục và phân hoạch dữ liệu nằm tại \path{experiments/training/results_mem_2/mem-adapter-v1}.
Mã băm của trọng số nền và HDF5 gốc được tham chiếu theo bản kê.

Mô phỏng chạy trong WSL với Python 3.8.20, robosuite 1.4.0, MuJoCo 3.2.3 và
NumPy 1.22.4. Suy luận trên Windows dùng Python 3.12.2, torch 2.6.0+cu124,
LeRobot 0.6.1, transformers 5.5.4 và PEFT 0.19.1.
Huấn luyện dùng torch 2.11.0+cu130, LeRobot 0.6.1, transformers 5.5.4,
PEFT 0.21.2 và accelerate 1.15.0. Phiên bản mô hình SmolVLA và nguồn tác vụ
LIBERO-Mem được ghi trong danh mục tài liệu tham khảo và trong bản kê từng lần chạy.

\section{Tệp dữ liệu phân tích và bối cảnh V1}
\path{docs/thesis/evidence/mem_test_episodes.csv} giữ một dòng cho mỗi
lượt, gồm kết quả, mã băm trạng thái, seed, tiến độ do bộ chấm và bộ ghi
xác định, cờ lịch sử và nguồn bản tóm tắt. \path{mem_test_counts.csv} tổng hợp
theo tác vụ. \path{validation_loss.csv} lấy từ nhật ký xác thực.
\path{historical_run_inventory.csv} là danh mục các lần chạy lịch sử.

Ở giai đoạn V1, báo cáo nội bộ ghi 77/160 lượt thành công trên 40 tác vụ,
28/40 trên tập nghiệm thu, và 11/40 trên LIBERO-10. Trong 83 lượt thất bại,
71 lượt chưa được quy nguyên nhân gốc theo tiêu chí chẩn đoán của báo cáo.
Tại mốc tổng hợp, các bản ghi từng lượt vắng mặt ở đường dẫn báo cáo dẫn chiếu.
Các số liệu này được dẫn theo báo cáo V1 và dùng làm bối cảnh lịch sử.
V1 dùng $s=50$ giữa hai lần truy vấn, còn đánh giá T1/T3 dùng $s=10$.
Khác biệt về nhịp điều khiển và tập tác vụ giới hạn khả năng đối chiếu giữa hai giai đoạn.

\section{Đối chiếu bản ghi sự kiện thị giác}
Phép đo ngày 26/10 được tổng hợp tại
\path{experiments/audits/memory_writer_v3/verification_summary.json}.
Nhãn tham chiếu đánh dấu mốc tiến độ tích lũy trên các quỹ đạo trình diễn.
Các báo cáo ghép mốc, kết quả thành phần và bản lưu mã nguồn được đối chiếu
SHA-256. Tập bằng chứng cũng ghi kết quả so khớp phiên bản của bộ theo dõi,
bộ ghi và bộ dựng văn bản với phiên bản dùng trong lần đo.

Các phép kiểm kỹ thuật đi kèm gồm đối chiếu đầu ra nhận dạng trên 12 khung hình,
đối chiếu chỉ mục trên 162 mẫu, 49 kiểm thử thành phần và 47 kiểm thử hồi quy.
Những phép kiểm này hỗ trợ đánh giá tính nhất quán của quy trình xử lý.
Độ phù hợp của sự kiện được phân tích bằng mốc chú giải và các quỹ đạo
trong Chương 4.

\section{Tái tổng hợp số liệu}
\path{docs/thesis/README.md} ghi các lệnh đọc bản ghi đã lưu, tổng hợp số liệu
và tái tạo tài liệu. Mỗi lần tổng hợp sử dụng tập bằng chứng và mã băm tương ứng.
